% Part I, introduction and related work (included by main_en.tex).

\section{Introduction}\label{sec:intro}

Reinforcement learning needs a reward, and for most tasks someone has to write it. The MuJoCo locomotion benchmarks
come with engineered rewards: forward velocity plus a survival bonus minus a control cost, with weights that were tuned
by hand. A language model that judges behaviour from a description of the task would remove that step. This idea has
been pursued in several forms. Language models have written reward code~\cite{yu2023language,ma2024eureka}, and they have
served directly as reward functions in text environments~\cite{kwon2023reward}. Vision--language models have scored
images of a robot's pose~\cite{rocamonde2024vlm}. And a language model's preferences over event captions have been
distilled into an intrinsic reward~\cite{klissarov2024motif}.

This paper asks a narrow version of the question, with the controls needed to answer it. Suppose a 4-billion-parameter
judge reads a short text summary of what the robot did in a window of 32 steps. Can its score replace the engineered
reward of a standard benchmark when PPO is trained on it, and what does that cost? We compare two judges of the same
size and family. One is \kev{}~\cite{kev4b}, a decision model that answers questions about a JSON state in a single
forward pass, with a pointer head and a calibration temperature. The other is an autoregressive (AR) judge,
Qwen3.5-4B, whose next-token distribution over the answer options is read out the same way. Both judges see identical
text. A control with the true reward at the judges' temporal resolution (\texttt{env\_window}) separates the effect of
window-level feedback from the effect of the judge.

\paragraph{Contributions.}
\begin{itemize}
\item A controlled comparison of seven reward sources for PPO on five MuJoCo tasks: the engineered reward, its
  window-averaged version, \kev{}'s score with and without a ``healthy'' gate, a distilled reward model labelled by
  \kev{} on 10\% of the windows, and the AR judge with expected and with sampled readout. It uses 1M steps and 5 seeds
  per task (3 for the sampled AR judge), and statistics after \citet{agarwal2021deep} (\S\ref{sec:results}).
\item An offline agreement study that relates each judge signal to the environment's reward on the same windows
  (\S\ref{sec:agreement}).
\item Measured costs: judged windows, judge time and wall-clock per run (\S\ref{sec:cost}).
\item A correction. Earlier versions of this paper (v1--v3) reported results that were not produced by these
  experiments. During training the ``judge'' was a hand-written function of the simulator state, and the leaderboards
  were constants (see the notice above and \texttt{CORRECTION.md}). Every number in this version is computed from the
  logs in the repository by the scripts there.
\end{itemize}

\section{Related work}\label{sec:related}

\paragraph{Language models as reward designers and reward functions.} Eureka~\cite{ma2024eureka} and Language to
Rewards~\cite{yu2023language} use a language model to write or parameterise reward code. This keeps the policy's reward
cheap but moves the design effort into the model's coding ability and into an outer search loop. Using the model as the
reward itself, as \citet{kwon2023reward} do for negotiation and other text games, needs no code, but the model is then
called during training. VLM-RM~\cite{rocamonde2024vlm} scores rendered frames with CLIP for humanoid poses. Motif~\cite{klissarov2024motif}
distils a language model's preferences between NetHack captions into a reward network, which our \texttt{kev\_distill}
condition resembles. Our judges read text summaries of telemetry instead of pixels. This isolates the judgment from
perception, at the price of a summary that a person had to design.

\paragraph{LLM judges.} Using language models to grade outputs~\cite{zheng2023judging,liu2023geval} has become standard
in evaluation. As reward models they are known to be exploitable~\cite{gao2023scaling,skalse2022defining}. Generative
verifiers~\cite{zhang2024genverifier} read a judge's yes/no probability, as our AR judge does. \kev{} differs in that its
answers are distributions from a dedicated head, calibrated with a temperature, and computed in one pass for several
questions.

\paragraph{Evaluation of deep RL.} With few seeds per task, point estimates of mean return are unreliable. We follow
\citet{agarwal2021deep}: interquartile means over all runs with stratified-bootstrap confidence intervals, and the
probability of improvement. We report per-task returns as well.
